\subsection{本原理想}

设 A 是 K 上的一个代数，E 是 K 上的一个向量空间。A 在 E 中的一个表示是指从 A 到 E 的自同态所成的代数 \(\mathcal{L}(\mathrm{E})\) 的一个态射。单射的表示称为忠实的。设 \(\pi_1\) 与 \(\pi_2\) 是 A 在空间 \(\mathrm{E}_1\)、\(\mathrm{E}_2\) 中的表示。从 \(\pi_1\) 到 \(\pi_2\) 的一个态射是指一个 K-线性映射 \(u: \mathrm{E}_1 \to \mathrm{E}_2\)，它使得 \(u(\pi_1(a)x) = \pi_2(a)u(x)\) 对一切 \(a \in \mathrm{A}\) 与 \(x \in \mathrm{E}_1\) 成立。若存在一个从 \(\pi_1\) 到 \(\pi_2\) 的态射，它是向量空间的一个同构，则称这些表示是等价的。其逆于是是从 \(\pi_2\) 到 \(\pi_1\) 的一个态射。A 在 E 中的表示 \(\pi\) 称为不可约的，如果 \(\mathrm{E} \neq \{0\}\)，并且 E 的在 \(\pi(\mathrm{A})\) 下稳定的向量子空间只有 \(\{0\}\) 与 E。

例. --- 从 A 到 \(\mathcal{L}(\mathrm{E})\) 的零映射是一个表示，称为 A 的平凡表示。它不可约必须且只需 E 的维数为 1。

引理 4. --- 设 π 是 A 在 E 中的一个不可约的非平凡表示。对 E 的每个非零元素 ξ，有 π(A)ξ = E。

子空间 \(\pi(A)\xi\)（\(E\) 的子空间）在 \(\pi(A)\) 下稳定。假定它为零。那么非零子空间 \(K\xi\)（\(E\) 的非零子空间）就会在 \(\pi(A)\) 下稳定，从而等于 \(E\)；但这将蕴涵 \(\pi\) 是零表示。因此我们有 \(\pi(A)\xi = E\)。

设 \(\pi\) 是 A 在 E 中的一个不可约的非平凡表示。根据该引理，\(\xi\) 在 A 中的零化子 R 是 A 的一个正则左理想 (A, VIII, p.~425, 第 1 小节)，且表示 \(\pi\) 等价于由 A-伪模 A/R 所定义的表示。由于 \(\pi\) 不可约，理想 R 是一个极大正则左理想。

定义 7. --- 设 A 是 K 上的一个代数。A 的本原理想是指 A 的一个非平凡不可约表示的核。

若 A 交换，则 A 的本原理想就是 A 的正则极大理想。事实上，A 的非平凡不可约表示，除了相差一个等价外，就是由 A-伪模 A/R 所定义的表示 \(\pi_{R}\)，其中 R 是 A 的一个正则极大理想。\(\pi_{R}\) 的核包含 R。它甚至等于 R，因为由 A, VIII, p.~426, 命题 2，A 的交换性蕴涵 A/R 是一个域。因此 \(\text{Ker}(\pi_{\text{R}})\) 是正则极大的。

引理 5. --- 设 π 是 A 在 K 上一个向量空间 E 中的一个不可约表示。

\begin{enumerate}
\def\labelenumi{\alph{enumi})}
\item
  设 I 是 A 的一个双边理想。若 \(\pi(\mathrm{I}) \neq \{0\}\)，则 \(\pi|\mathrm{I}\) 不可约；
\item
  设 \(I_{1}\) 与 \(I_{2}\) 是 A 的双边理想，使得 \(\pi(\mathrm{I}_{1}) \neq 0\) 且 \(\pi(\mathrm{I}_{2}) \neq 0\)。则 \(\pi(\mathrm{I}_{1}\mathrm{I}_{2}) \neq 0\)。
\end{enumerate}

E 中被 \(\pi(\mathrm{I})\) 零化的元素所成的集合在 \(\pi(\mathrm{A})\) 下稳定且异于 E，故等于 0。于是，若 \(\xi\) 是 E 的一个非零元素，则有 \(\pi(\mathrm{I})\xi \neq 0\)；由于 \(\pi(\mathrm{I})\xi\) 在 \(\pi(\mathrm{A})\) 下稳定，有 \(\pi(\mathrm{I})\xi = \mathrm{E}\)，这就证明了 \(a\) )。另一方面，上文证明了 \(\pi(\mathrm{I}_2)\mathrm{E} = \mathrm{E}\)，\(\pi(\mathrm{I}_1)\pi(\mathrm{I}_2)\mathrm{E} = \mathrm{E}\)，从而 \(\pi(\mathrm{I}_1\mathrm{I}_2) \neq 0\)，由此得 \(b\) )。

引理 6. --- 设 I \(_{1}\) 与 I \(_{2}\) 是 A 的双边理想，I 是 A 的一个本原理想。若 I 包含 I \(_{1}\) I \(_{2}\)（特别地，若 I 包含 I \(_{1}\) ∩ I \(_{2}\)），则 I 包含 I \(_{1}\) 或 I \(_{2}\)。

设 \(\pi\) 是以 I 为核的一个不可约表示。若 I \(\not\supseteq\) I \(_1\) 且 I \(\not\supseteq\) I \(_2\)，则引理 5 的 \(b\) ) 证明 \(\pi(\text{I}_1\text{I}_2) \neq 0\)，从而 I \(\not\supseteq\) I \(_1\) I \(_2\)。

引理 7. --- 假定 A 有单位元。设 I 是 A 的一个极大双边理想。则 I 是一个本原理想。

存在 A 的一个包含 I 的极大左理想 R (A, I, p.~99, 定理 1)。设 \(\pi\) 是 A 在 A/R 中的典范表示，它不可约且非零。由于 IA \(\subset\) R，\(\pi\) 的核 I′ 包含 I，故 I′ = I，从而 I 是本原的。

设 J(A) 是 A 的本原理想所成的集合。对 A 的任何子集 M，我们用 V(M) 表示包含 M 的 A 的本原理想所成的集合；若 I 是由 M 生成的 A 的双边理想，则 V(M) = V(I)。若 M 化为单个元素 x，我们就写 V(x) 以代替 V(\{x\})。

映射 M \(\mapsto\) V(M) 关于包含关系是递减的。我们有：

\[
\begin{array}{r l} (5) & \mathrm{V} (\varnothing) = \mathrm{J} (\mathrm{A}), \qquad \mathrm{V} (\mathrm{A}) = \varnothing \end{array}
\]

\[
\begin{array}{r l} (6) & \mathrm{V} \left(\bigcup_ {i \in \mathrm{I}} \mathrm{M} _ {i}\right) = \mathrm{V} \left(\sum_ {i \in \mathrm{I}} \mathrm{M} _ {i}\right) = \bigcap_ {i \in \mathrm{I}} \mathrm{V} (\mathrm{M} _ {i}) \end{array}
\]

对 A 的子集的任何族 \((\mathrm{M}_{i})_{i\in\mathrm{I}}\) 成立。另一方面，由引理 6，

\[
\begin{array}{r l} (7) & \mathrm{V} (\mathrm{I} _ {1} \cap \mathrm{I} _ {2}) = \mathrm{V} (\mathrm{I} _ {1} \mathrm{I} _ {2}) = \mathrm{V} (\mathrm{I} _ {1}) \cup \mathrm{V} (\mathrm{I} _ {2}) \end{array}
\]

对 A 的任意双边理想 \(I_{1}\)、\(I_{2}\) 成立。公式 (5) 至 (7) 表明，J(A) 的子集 V(M) 是一个称为 J(A) 上 Jacobson 拓扑的拓扑的闭子集。

设 T 是 J(A) 的一个子集，并设 \(\Upsilon(T)\) 是 T 的诸元素之交，于是 \(\Upsilon(T)\) 是 A 的一个双边理想。则 T 在 J(A) 中的闭包是包含 T 的最小闭子集，即 V( \(\Upsilon(T)\) )。特别地，T 是闭的必须且只需 \(T = V(\Upsilon(T))\)。

命题 2. --- 设 I \(_{1}\) 与 I \(_{2}\) 是 J(A) 的两个不同的点。则这两点中的一个不附着于另一个。

事实上，例如我们有 \(I_{1} \not\subset I_{2}\)。集合 \(V(I_{1})\)，它由 \(I \in J(A)\)（使得 \(I_{1} \subset I\)）组成，在 \(J(A)\) 中是闭的，它包含 \(I_{1}\) 但不包含 \(I_{2}\)。

命题 3. --- 设 I ∈ J(A)。要使 \{I\} 在 J(A) 中是闭的，必须且只需 I 是一个本原极大理想。

事实上，\{I\} 的闭包由包含 I 的 A 的本原理想组成。

关系

« \(\pi_{1},\pi_{2}\) 是 A 的表示，且它们同构。

在 \(\pi_1\) 与 \(\pi_2\) 之间是一个等价关系。对 A 的每个表示 \(\pi\)，我们用 \(\operatorname{cl}(\pi)\) 表示 \(\pi\) 的等价类，它因此是一个与 \(\pi\) 同构的 A 的表示，使得两个表示 \(\pi_1\) 与 \(\pi_2\) 同构必须且只需 \(\operatorname{cl}(\pi_1) = \operatorname{cl}(\pi_2)\)。我们说 \(\operatorname{cl}(\pi)\) 是 \(\pi\) 的类。

设 c 是 A 的基数。设 \(\pi\) 是 A 在一个 K-向量空间 E 中的一个非零不可约表示。设 \(\xi\) 是 E 的一个非零元素。由于 \(\pi(A)\xi = E\) (引理 4)，E 的维数 \(\leqslant\mathfrak{c}\) (A, II, p.~97, 推论)。关系

“λ 是 A 的不可约表示的一个类

在一个维数 \(\leqslant\mathfrak{c}\gg\)

在 \(\lambda\) 中是可汇集的 (E, II, p.~3)。事实上，每个维数 \(\leqslant \mathfrak{c}\) 的向量空间都同构于一个空间 \(\mathrm{K}^{\mathrm{B}}\)，其中 B 是 A 的一个子集 (A, II, p.~25, 定义 10)，而断言则由 E, II, p.~47 得出。

我们用 \(\widehat{\mathrm{A}}\) 表示 A 的非平凡不可约表示的类所成的集合。根据上文，对 A 的每个非平凡不可约表示 \(\pi\)，存在唯一的一个表示 \(\widehat{\pi} \in \widehat{\mathrm{A}}\)，它与 \(\pi\) 同构。

从 \(\hat{\mathrm{A}}\) 到 \(\mathrm{J(A)}\) 中的、把 \(\pi\) 对应到其核的映射是满射。若 A 交换，则由本原理想是正则极大理想这一事实推出，这个映射是一个双射。

我们给 \(\hat{A}\) 赋予在映射 \(\hat{A} \rightarrow J(A)\) 下的原像拓扑。

若 A 有单位元，则空间 J(A) 与 \(\hat{\mathrm{A}}\) 是拟紧的。

只须对 J(A) 给出证明。设 \((\mathrm{T}_{j})\) 是 J(A) 的闭子集的一个族，其交为空。若和 \(\sum_{j}\Upsilon(\mathrm{T}_{j})\) 异于 A，则这个和将包含于一个极大双边理想 I。理想 I 将是本原的 (引理 7)；由于每个 \(T_{j}\) 是闭的，故等于 \(\mathrm{V}(\Upsilon(\mathrm{T}_{j}))\)，我们就会有 \(I\in T_{j}\) 对一切 j 成立，这与假设矛盾。于是我们有 \(\sum_{j}\Upsilon(\mathrm{T}_{j})=\mathrm{A}\)，因而可以写 \(1=x_{1}+\cdots+x_{n}\)，其中 \(n\geqslant1\)，且对一切 i 有 \(x_{i}\in\Upsilon(\mathrm{T}_{j_{i}})\)。这蕴涵 \(\Upsilon(\mathrm{T}_{j_{1}})+\cdots+\Upsilon(\mathrm{T}_{j_{n}})=\mathrm{A}\)，从而 \(T_{j_{1}}\cap\cdots\cap T_{j_{n}}=\varnothing\)。

假定代数 A 是交换的且有单位元。J(A) 上的 Jacobson 拓扑是由 A 的素谱的 Zariski 拓扑在 J(A) 上诱导的拓扑 (AC, II, 定义 4, p.~125)。

假定 A 交换且 K 是一个拓扑域。从 K 到 \(\mathcal{L}(\mathrm{K})\) 上的典范同构使我们可以把 \(\mathsf{X}(\mathsf{A})\) 的一个元素等同于 A 在向量空间 K 中的一个表示，这就定义了一个从 \(\mathsf{X}(\mathsf{A})\) 到 \(\widehat{\mathsf{A}}\) 中的单射。因此可以把 \(\mathsf{X}(\mathsf{A})\) 与 \(\widehat{\mathsf{A}}\) 的一个子集视为同一。

命题 5. — 在 \(\mathsf{X}(\mathsf{A})\) 上由 \(\widehat{\mathsf{A}}\) 的拓扑诱导的拓扑粗于 \(\mathsf{X}(\mathsf{A})\) 的拓扑。

事实上，设 T 是 \(\widehat{\mathbf{A}}\) 的一个闭子集。则 T 是由其核包含 A 的一个子集 M 的那些 \(\pi \in \widehat{\mathbf{A}}\) 所成的集合。因此 \(\mathrm{T} \cap \mathrm{X}(\mathrm{A})\) 是由在 M 上为零的那些 \(\chi \in \mathrm{X}(\mathrm{A})\) 所成的集合，即 \(\mathrm{X}(\mathrm{A})\) 的一个闭子集。由此得命题。

一般说来，\(\mathsf{X}(\mathsf{A})\) 的拓扑与由 \(\widehat{\mathsf{A}}\) 的拓扑诱导的拓扑并不重合 (参见 I, p.~193, 习题 6, c))。

